\chapter{Vergleich und wiederholte Quotientenbildung}
\label{b11guide:ch:quotients}

\section{Eine Äquivalenzrelation ist in einer anderen enthalten}

\begin{theorem}[Inklusion, Klassen und gesättigte Mengen]
\label{b11guide:thm:inclusion}
Für $E,F\in\EqSet(A)$ sind äquivalent:
\[
\begin{aligned}
 &E\subseteq F,\\
 &\forall x\in A\;([x]_E\subseteq[x]_F),\\
 &\forall C\in A/E\;\exists D\in A/F\;(C\subseteq D),\\
 &\text{jede }F\text{-gesättigte Menge ist }E\text{-gesättigt}.
\end{aligned}
\]
Falls diese Bedingungen gelten, ist jeder $F$-Block
die Vereinigung genau der in ihm enthaltenen $E$-Blöcke:
\[
 D=\bigcup\{C\in A/E\mid C\subseteq D\}\qquad(D\in A/F).
\]
\end{theorem}

\begin{theorem}[Sättigung bei ineinander enthaltenen Relationen]
Seien $E,F\in\EqSet(A)$ mit $E\subseteq F$.
Für $B\subseteq A$ gilt
\[
\begin{aligned}
 \Sat_E(B)&\subseteq\Sat_F(B),\\
 \Sat_F(\Sat_E(B))&=\Sat_F(B),\\
 \Sat_E(\Sat_F(B))&=\Sat_F(B).
\end{aligned}
\]
\end{theorem}

\section{Eine Relation auf der Menge der Klassen}

\begin{definition}[Induzierte Äquivalenzrelation auf dem Quotienten]
\label{b11guide:def:quotientrelation}
Seien $E,F\in\EqSet(A)$ und $E\subseteq F$. Auf $A/E$ sei
$F/E$ durch
\[
 C\sim_{F/E}D
 \quad\Longleftrightarrow\quad
 \exists x\in C\;\exists y\in D\;(x\sim_F y)
 \qquad(C,D\in A/E)
\]
definiert.
\end{definition}

\begin{theorem}[Unabhängigkeit von Repräsentanten]
\label{b11guide:thm:quotientrelation}
Für $E,F\in\EqSet(A)$ mit $E\subseteq F$ gilt
$F/E\in\EqSet(A/E)$. Für $x,y\in A$ gilt
\[
 [x]_E\sim_{F/E}[y]_E\quad\Longleftrightarrow\quad x\sim_F y.
\]
Für Klassen $C,D\in A/E$ sind äquivalent:
\[
 C\sim_{F/E}D,\qquad
 \exists x\in C\;\exists y\in D\;(x\sim_F y),\qquad
 \forall x\in C\;\forall y\in D\;(x\sim_F y).
\]
\end{theorem}

\begin{theorem}[Klassen des wiederholten Quotienten]
\label{b11guide:thm:iterated}
Für $E,F\in\EqSet(A)$ mit $E\subseteq F$ und $x\in A$ gilt
\[
\begin{aligned}
 [[x]_E]_{F/E}
   &=\{C\in A/E\mid C\subseteq[x]_F\},\\
 \bigcup [[x]_E]_{F/E}&=[x]_F.
\end{aligned}
\]
Damit gilt die Mengenform des Satzes über wiederholte Quotienten:
\[
 A/F=\left\{\bigcup\mathcal D\;\middle|\;
                     \mathcal D\in(A/E)/(F/E)\right\}.
\]
Verschiedene Klassenfamilien $\mathcal D$ des wiederholten
Quotienten haben verschiedene Vereinigungen.
Die beiden Quotienten sind im allgemeinen keine gleichen Mengen:
Ein Element von $(A/E)/(F/E)$ ist eine Familie von $E$-Klassen,
ein Element von $A/F$ eine Teilmenge von $A$.
\end{theorem}

\Needspace{19\baselineskip}
\section{Alle Äquivalenzrelationen auf einem Quotienten}

\begin{definition}[Aus Klassen zurückgewonnene Relation]
Seien $E\in\EqSet(A)$ und $T\in\EqSet(A/E)$.
Die Relation $\widetilde T$ auf $A$ sei definiert durch
\[
 x\sim_{\widetilde T}y
 \quad\Longleftrightarrow\quad [x]_E\sim_T[y]_E.
\]
\end{definition}

\begin{theorem}[Entsprechung der Äquivalenzrelationen]
\label{b11guide:thm:correspondence}
Seien $E\in\EqSet(A)$ und $T\in\EqSet(A/E)$. Dann gilt
\[
 \widetilde T\in\EqSet(A),\qquad
 E\subseteq\widetilde T,\qquad \widetilde T/E=T.
\]
Umgekehrt gilt für jedes $F\in\EqSet(A)$ mit $E\subseteq F$
\[
 \widetilde{(F/E)}=F.
\]
Jede Äquivalenzrelation auf $A/E$ wird damit von genau einer
Äquivalenzrelation auf $A$, die $E$ enthält, induziert.
Ihre Klassen erfüllen
\[
 [x]_{\widetilde T}=\bigcup [[x]_E]_T\qquad(x\in A).
\]
Die Aussagen umfassen $A=\varnothing$.
\end{theorem}

